\section{Limitations and broader impact}

\sys is a bounded benchmark, not a claim that commodity hardware can substitute for production-scale evaluation. The experiments run on one local operator machine, so the results do not characterize distributed retrieval, GPU indexing, cloud variability, or multi-tenant production workloads. S3 uses \textsc{HotpotQA-portable}, a frozen paragraph-level proxy for full-Wikipedia multi-hop retrieval, and S2 measures top-k retrievability after inserts rather than lower-level index visibility. These choices make the benchmark reproducible and one-day executable, but they narrow the scope of the claims.

The uncertainty estimates are conditional on this fixed benchmark setting. Confidence intervals are query- or event-level summaries, not hardware-population intervals, and observed p99 latency is not a production tail-latency guarantee. The current result bundle also shows why short benchmark runs should not be overinterpreted. S1 and S2 change top-1 recommendations between B0 and B2, so a screening run is not enough to make a final deployment decision for those workloads. S3 keeps the same top-1 recommendation but has low full-rank correlation, indicating that the best decision can be stable even when the rest of the ranking is not. Future versions should add more repeated B2 runs, more event streams, dedicated S2 competitor reruns, a second hardware profile, and an explicit production-sized stress track.

The broader impact of \sys is primarily methodological. A reproducible benchmark can help operators choose lower-cost, lower-latency retrieval infrastructure without hiding quality or freshness tradeoffs. The negative risk is that readers may treat the benchmark winner as a universal recommendation rather than a result under a specific machine, dataset, objective, and conformance policy. The manuscript mitigates that risk by reporting exclusions, limitations, budget stability, and objective sensitivity directly.
